% GENERATED FILE. Edit canonical lessons, then run npm run generate:books.
% canonical-source-hash: fnv1a64:cd1d9f456160ee69
% canonical-lessons: RU-C138-budushchiy, RU-C138-cherez, RU-C138-plan, RU-C138-planirovat, RU-C138-nadeyatsya

\chapter{Future, in a Week, a Plan, To Plan, To Hope}
\label{ch:ru-future-in-a-week-a-plan-to-plan-to-hope}
\hlchaptermodality{138}

\begin{chapteropening}
\textbf{By the end of this chapter:} \emph{I can say what will happen, and when: next year, in a week, and what I plan.}

Complete the last lesson of chapter 138: I can say what will happen, and when: next year, in a week, and what I plan.
\end{chapteropening}

\section[\texorpdfstring{\ru{будущий} (búdushchiy)}{будущий (búdushchiy)}]{\ru{будущий} (búdushchiy) — future, next}

\label{lesson:RU-C138-budushchiy}



\begin{quote}
Before the new one: say the Russian for last, then the Russian for to happen.
\end{quote}

\subsection*{You'll want to know: \ru{будущий}}
\textbf{\ru{будущий}} (\emph{búdushchiy}) — \textquotedblleft{}future, next\textquotedblright{}. One Russian future is \textbf{\ru{буду}} and an infinitive: \textbf{\ru{Я} \ru{буду} \ru{читать}.} — \emph{ya búdu chitát} — \textquotedblleft{}I will read.\textquotedblright{} \textbf{\ru{в} \ru{будущем} \ru{году}} is \textquotedblleft{}next year\textquotedblright{}.

\begin{grammarlens}[title={What you've built}]
One more word for this chapter.
\end{grammarlens}

\subsection*{Your turn}
\begin{itemize}
\raggedright
  \item \emph{Say it:} \emph{búdushchiy}
  \item \emph{Say it:} \emph{búdushchiy}, once more
  \item \emph{Say it:} say \emph{búdushchiy}
  \item \emph{Recall:} say the Russian for last, then the Russian for to happen, then say \emph{búdushchiy} again
\end{itemize}

\begin{tcolorbox}[breakable,colback=teal!4,colframe=teal!35!black,title={Before you move on}]
What does \ru{будущий} mean? (\textquotedblleft{}Future, next\textquotedblright{}.) Say it once more.
\end{tcolorbox}
\section[\texorpdfstring{\ru{через} (chérez)}{через (chérez)}]{\ru{через} (chérez) — in (a week), after}

\label{lesson:RU-C138-cherez}



\begin{quote}
Before the new one: say the Russian for to happen, then the Russian for future.
\end{quote}

\subsection*{You'll want to know: \ru{через}}
\textbf{\ru{через}} (\emph{chérez}) — \textquotedblleft{}in (a week), after\textquotedblright{}. Counted from now: \textbf{\ru{через} \ru{неделю}} — \emph{chérez nedélyu} — \textquotedblleft{}in a week\textquotedblright{}. \textbf{\ru{неделя}} changes to \textbf{\ru{неделю}} after it.

\begin{grammarlens}[title={What you've built}]
One more word for this chapter.
\end{grammarlens}

\subsection*{Your turn}
\begin{itemize}
\raggedright
  \item \emph{Say it:} \emph{chérez}
  \item \emph{Say it:} \emph{chérez}, once more
  \item \emph{Say it:} say \emph{chérez}
  \item \emph{Recall:} say the Russian for to happen, then the Russian for future, then say \emph{chérez} again
\end{itemize}

\begin{tcolorbox}[breakable,colback=teal!4,colframe=teal!35!black,title={Before you move on}]
What does \ru{через} mean? (\textquotedblleft{}In (a week), after\textquotedblright{}.) Say it once more.
\end{tcolorbox}
\section[\texorpdfstring{\ru{план} (plan)}{план (plan)}]{\ru{план} (plan) — a plan}

\label{lesson:RU-C138-plan}



\begin{quote}
Before the new one: say the Russian for future, then the Russian for in.
\end{quote}

\subsection*{You'll want to know: \ru{план}}
\textbf{\ru{план}} (\emph{plan}) — \textquotedblleft{}a plan\textquotedblright{}. \textbf{\ru{Какой} \ru{план}?} — \emph{kakóy plan} — \textquotedblleft{}What's the plan?\textquotedblright{} English \emph{plan} is the same word.

\begin{grammarlens}[title={What you've built}]
One more word for this chapter.
\end{grammarlens}

\subsection*{Your turn}
\begin{itemize}
\raggedright
  \item \emph{Say it:} \emph{plan}
  \item \emph{Say it:} \emph{plan}, once more
  \item \emph{Say it:} say \emph{plan}
  \item \emph{Recall:} say the Russian for future, then the Russian for in, then say \emph{plan} again
\end{itemize}

\begin{tcolorbox}[breakable,colback=teal!4,colframe=teal!35!black,title={Before you move on}]
What does \ru{план} mean? (\textquotedblleft{}A plan\textquotedblright{}.) Say it once more.
\end{tcolorbox}
\section[\texorpdfstring{\ru{планировать} (planírovat)}{планировать (planírovat)}]{\ru{планировать} (planírovat) — to plan}

\label{lesson:RU-C138-planirovat}



\begin{quote}
Before the new one: say the Russian for in, then the Russian for a plan.
\end{quote}

\subsection*{You'll want to know: \ru{планировать}}
\textbf{\ru{планировать}} (\emph{planírovat}) — \textquotedblleft{}to plan\textquotedblright{}. \textbf{\ru{Я} \ru{планирую} \ru{путешествие}.} — \emph{ya planíruyu puteshéstviye} — \textquotedblleft{}I'm planning a trip.\textquotedblright{} Verbs in \textbf{-\ru{овать}} swap it for \textbf{-\ru{ую}}, like \textbf{\ru{рисовать}}, \textbf{\ru{рисую}}.

\begin{grammarlens}[title={What you've built}]
One more word for this chapter.
\end{grammarlens}

\subsection*{Your turn}
\begin{itemize}
\raggedright
  \item \emph{Say it:} \emph{planírovat}
  \item \emph{Say it:} \emph{planírovat}, once more
  \item \emph{Say it:} say \emph{planírovat}
  \item \emph{Recall:} say the Russian for in, then the Russian for a plan, then say \emph{planírovat} again
\end{itemize}

\begin{tcolorbox}[breakable,colback=teal!4,colframe=teal!35!black,title={Before you move on}]
What does \ru{планировать} mean? (\textquotedblleft{}To plan\textquotedblright{}.) Say it once more.
\end{tcolorbox}
\section[\texorpdfstring{\ru{надеяться} (nadéyatsya)}{надеяться (nadéyatsya)}]{\ru{надеяться} (nadéyatsya) — to hope}

\label{lesson:RU-C138-nadeyatsya}



\begin{quote}
Before the new one: say the Russian for a plan, then the Russian for to plan.
\end{quote}

\subsection*{You'll want to know: \ru{надеяться}}
\textbf{\ru{надеяться}} (\emph{nadéyatsya}) — \textquotedblleft{}to hope\textquotedblright{}. \textbf{\ru{Я} \ru{надеюсь}, \ru{что} ...} — \emph{ya nadéyus, shto} — \textquotedblleft{}I hope that ...\textquotedblright{} Like \textbf{\ru{подниматься}}, it keeps its \textbf{-\ru{ся}}, written \textbf{-\ru{сь}} after a vowel: \textbf{\ru{надеюсь}}.

\begin{grammarlens}[title={What you've built}]
One more word for this chapter.
\end{grammarlens}

\subsection*{Your turn}
\begin{itemize}
\raggedright
  \item \emph{Say it:} \emph{nadéyatsya}
  \item \emph{Say it:} \emph{nadéyatsya}, once more
  \item \emph{Say it:} say \emph{nadéyatsya}
  \item \emph{Recall:} say the Russian for a plan, then the Russian for to plan, then say \emph{nadéyatsya} again
\end{itemize}

\begin{tcolorbox}[breakable,colback=teal!4,colframe=teal!35!black,title={Before you move on}]
What does \ru{надеяться} mean? (\textquotedblleft{}To hope\textquotedblright{}.) Say it once more.
\end{tcolorbox}
